\typeout{NT FILE chapter-integration-testing.tex}%
\chapter{Integration and Testing}
\label{cha:integration-testing}

This chapter describes how the components developed during this internship, the
\gls{sdr-control-system} and the \gls{rs232-interface}, were brought together
into a working system, the environment in which that system was validated, and
the testing strategy used to ensure its correctness and reliability.
\section{Unit Testing}

Before integrating the components, unit tests were developed for both the
\gls{sdr-control-system} and the \gls{rs232-interface}. These tests were
designed to validate the functionality of individual components in isolation,
ensuring that each component behaved as expected under a variety of conditions.
For the \gls{sdr-control-system}, unit tests were created to verify the correct
generation of \gls{GNSS} signals, the proper handling of configuration
parameters, and the correct interaction with the underlying \gls{SDR} hardware.
For the \gls{rs232-interface}, unit tests deemed to be harder to implement as
the hardware used did not have same port loopback capabilities, thus
unfortunately only framing, serialization, and protocol dispatch were unit
tested. Proper \gls{RS232} testing required the development of the external
tester described in \autoref{sec:the-external-tester}.
\section{Environment Overview}

Throughout this chapter, references to ``the computer'', ``the host'' or ``the
generator'' refer specifically to the machine running the \gls{GNSS} signal
generator software stack, which hosts both the \gls{sdr-control-system} and the
\gls{rs232-interface}, and to which the \gls{SDR} hardware (DTA-2116 PCIe
cards) is physically attached.
It is important to frame the environment in which this system was developed and
validated as a prototype setup, distinct from a final deployment environment.
In a deployed scenario, the generator would be placed in a loopback, over a
\gls{RF} uplink to a remote receiver that would measure the transmitted signal,
and calculate any needed adjustments to the signal generator's configuration.
That remote receiver would then communicate those adjustments back to the
signal generator over the \gls{RS232} interface, closing the loop and allowing
the signal generator to adapt its transmission parameters in response to the
measured signal quality. In this prototype environment, however, such setup was
too complex and impractical to maintain for day-to-day development and testing,
and so a testing component was developed as described in
\autoref{sec:the-external-tester}, with the external tester component
simulating the behaviour of the remote receiver and any other measurement
equipment that would normally be present in a deployed scenario.
\section{The External Tester}
\label{sec:the-external-tester}
The external tester is a hardware/software component that simulates the
behaviour of the remote receiver and any other measurement equipment that would
normally be present in a deployed scenario. It is used to validate the
\gls{rs232-interface} and the \gls{sdr-control-system} in a controlled and
repeatable manner. The setup is quite simple, the generator is connected to
measurement equipment over a \gls{RF} link, which relays information to the
tester. The tester then communicates with the generator over a \gls{RS232}
link, issuing commands and producing responses that are representative of what
the real equipment would return. This allows the both the \gls{rs232-interface}
and the \gls{sdr-control-system} to be exercised end-to-end, covering message
framing, serialisation, protocol dispatch, handler invocation, and signal
generation. Thus allowing the behaviour of the real system to be validated
under controlled and repeatable conditions, without depending on the
availability of physical instrumentation.
\section{Integration}
\label{sec:integration}

During integration with the existing proprietary systems at GMV, it
became apparent that the gRPC client-server boundary described in
\autoref{sec:sdr-design-and-implementation}, while functioning correctly and
as designed, was not exercised in the final integration. The
proprietary system integrating with the \gls{sdr-control-system}
links the \texttt{app} component directly as a library and invokes
it in-process, rather than communicating with it through the gRPC
server. The gRPC server is implemented as a separate module within
the system, and in this integration context that module is simply
not built into the resulting binary, making the client-server hop
entirely absent from this deployment path rather than merely unused
at runtime.

This does not retroactively invalidate the architectural decision
described in \autoref{sec:sdr-design-and-implementation}: the gRPC layer
remains functional and independently testable through the CLI
client (\autoref{sec:sdr-functionality-and-use-cases}), and would be
exercised by any future integration that requires the
\gls{sdr-control-system} to be controlled remotely or from a
separate process. The finding instead reflects a mismatch between
the generality assumed at design time, driven by the original
goal of remote, multi-client control
(\autoref{sec:sdr-objective-and-motivation}), and the specific
integration constraints of the proprietary system actually
encountered, which favoured direct, in-process linkage over
out-of-process communication. This mirrors, at a smaller scale, the
broader generalisation lesson discussed in
\autoref{sec:design-patterns:sdr-iterations}.